\documentclass[12pt]{amsart}

\usepackage{amsmath,amssymb,amsthm}
\usepackage{booktabs}
\usepackage{hyperref}
\usepackage{enumitem}

\newtheorem{theorem}{Theorem}[section]
\newtheorem{proposition}[theorem]{Proposition}
\newtheorem{lemma}[theorem]{Lemma}
\newtheorem{corollary}[theorem]{Corollary}
\theoremstyle{definition}
\newtheorem{definition}[theorem]{Definition}
\newtheorem{conjecture}[theorem]{Conjecture}
\theoremstyle{remark}
\newtheorem{remark}[theorem]{Remark}
\newtheorem{question}[theorem]{Question}

\newcommand{\leg}[2]{\left(\frac{#1}{#2}\right)}

\title[Toward Closure of the Gateway Gap]{Toward Closure of the Gateway Gap:\\
NQR Targets, Subgroup Failure, and the Multi-$A$ Strategy}

\author{Macdara \'O Murch\'u}

\begin{document}

\begin{abstract}
We analyse the remaining open case of the gateway decomposition framework
for the Erd\H{o}s--Straus conjecture.
The key new observation is that for every prime $A \equiv 3 \pmod{4}$,
the target residue $t_A = -4N_A^2 \bmod A$
is always a quadratic non-residue (\textsc{NQR}) modulo $A$
(Lemma~\ref{lem:nqr-target}),
a consequence of $(-1/A) = -1$.
This implies that the gateway with parameter $A$ fails for prime $p$
\emph{if and only if} every prime factor of $N_A = (p+A)/4$
is a quadratic residue (\textsc{QR}) modulo $A$
(Theorem~\ref{thm:failure-char}).
We leverage this to characterise simultaneous failure across multiple $A$ values
as a joint \emph{$Q_A$-smoothness} condition on a family of related integers.
A conditional result (Theorem~\ref{thm:conditional}) is obtained
assuming a standard equidistribution hypothesis.
We identify the precise analytic statement
(Conjecture~\ref{conj:equidist}) that would complete the proof unconditionally,
and pose nine questions connecting this problem to
Galois theory, the large sieve, and the anatomy of integers.
\end{abstract}

\maketitle

\tableofcontents

% ============================================================
\section{Introduction}
% ============================================================

The companion paper~\cite{Gateway2026} introduces the \emph{gateway decomposition}:
for any prime $p$ and odd $A \equiv 3 \pmod{4}$ with $4 \mid (p+A)$,
a divisor $d \mid N_A^2$ (where $N_A = (p+A)/4$) satisfying
$d \equiv -p^2 \cdot 4^{-1} \pmod{A}$ yields an explicit representation
$4/p = 1/x + 1/y + 1/z$.
The full Erd\H{o}s--Straus conjecture reduces to showing that for every prime $p$
there exists at least one such pair $(A,d)$.
Algebraic arguments cover approximately $97.6\%$ of all primes;
computational verification covers the remainder through $p \le 10^9$
using $A \le 239$.

The residual open problem is:

\begin{quote}
\textbf{Open Problem.}
For every Case~B QR7 prime---i.e., a prime $p \equiv 1 \pmod{24}$ with
$p \equiv 1, 2,$ or $4 \pmod{7}$ and every prime factor of $(p+3)/4$
congruent to $1 \pmod{3}$---does there exist a prime $A \equiv 3 \pmod{4}$
with $4 \mid (p+A)$ and a divisor $d$ of $N_A^2 = \bigl((p+A)/4\bigr)^2$
satisfying $d \equiv -p^2 \cdot 4^{-1} \pmod{A}$?
\end{quote}

This paper provides three things: (1)~a clean algebraic characterisation of
gateway failure (Section~\ref{sec:nqr}--\ref{sec:failure});
(2)~a conditional theorem reducing the problem to a standard conjecture
(Section~\ref{sec:conditional});
and (3)~a precise statement of the analytic gap with nine targeted questions
(Sections~\ref{sec:gap}--\ref{sec:questions}).

% ============================================================
\section{The NQR Target Lemma}\label{sec:nqr}
% ============================================================

Throughout, $A$ is an odd prime with $A \equiv 3 \pmod{4}$,
$p$ a prime with $4 \mid (p+A)$, and $N_A = (p+A)/4$.
From Remark~2.5 of~\cite{Gateway2026}, the target residue is
\[
  t_A \;=\; -4N_A^2 \;\bmod A
  \;\equiv\; -p^2 \cdot 4^{-1} \;\pmod{A}.
\]

\begin{lemma}[NQR Target]\label{lem:nqr-target}
For every prime $A \equiv 3 \pmod{4}$ and every prime $p$ with $A \nmid N_A$,
the target $t_A$ is a \emph{quadratic non-residue} modulo~$A$.
\end{lemma}

\begin{proof}
\[
  \leg{t_A}{A}
  = \leg{-4N_A^2}{A}
  = \leg{-4}{A}\cdot\leg{N_A}{A}^2
  = \leg{-1}{A}\cdot\leg{4}{A}\cdot 1
  = \leg{-1}{A}.
\]
Since $A \equiv 3 \pmod{4}$, Euler's criterion gives $(-1/A) = -1$.
\end{proof}

\begin{remark}
The restriction $A \equiv 3 \pmod{4}$ is not a convenience but a structural
necessity: for $A \equiv 1 \pmod{4}$, the target would be a QR and the
failure characterisation below would break.
\end{remark}

% ============================================================
\section{Subgroup Failure Characterisation}\label{sec:failure}
% ============================================================

Let $G_A = (\mathbb{Z}/A\mathbb{Z})^\times$ (cyclic of order $A-1$)
and $Q_A \le G_A$ the unique index-2 subgroup (quadratic residues mod $A$).

\begin{definition}[$Q_A$-smooth]
A positive integer $n$ is \emph{$Q_A$-smooth} if every prime factor $q$ of $n$
satisfies $(q/A) = +1$ (i.e., $q \in Q_A$).
\end{definition}

\begin{theorem}[Failure Characterisation]\label{thm:failure-char}
Gateway $A$ fails for prime $p$ (i.e., no divisor of $N_A^2$
lies in the class $t_A \bmod A$) if and only if $N_A$ is $Q_A$-smooth.
\end{theorem}

\begin{proof}
The set of residues of divisors of $N_A^2$ modulo $A$ is
\[
  S_A \;=\;\Bigl\{\prod_{q \mid N_A} q^{b_q} \bmod A \;:\; 0 \le b_q \le 2v_q(N_A)\Bigr\}.
\]
If $N_A$ is $Q_A$-smooth, then every element of $S_A$ is a product of QRs,
hence a QR.  Since $t_A$ is a NQR (Lemma~\ref{lem:nqr-target}),
$t_A \notin S_A$ and the gateway fails.

Conversely, if $N_A$ has a prime factor $q$ with $(q/A) = -1$
then $q \in S_A$ (taking $b_q = 1$, $b_{q'} = 0$ for $q' \ne q$).
Since $q$ is a NQR, $S_A$ contains NQR elements.
Whether the \emph{specific} NQR target $t_A$ lies in $S_A$ depends on the full
factorisation of $N_A$ (and is handled by the character sum criterion
in Section~\ref{sec:charsum}); but $Q_A$-smoothness is the precise
obstruction to having \emph{any} NQR divisor at all.

More precisely: if $N_A$ is not $Q_A$-smooth, then $S_A$ contains both
QR and NQR elements (closed under multiplication within their respective cosets).
The gateway fails iff $t_A \notin S_A$, which (given $t_A$ is NQR)
requires $S_A$ to miss the specific coset-element $t_A$.
\end{proof}

\begin{remark}
Theorem~\ref{thm:failure-char} provides a sharp \emph{necessary} condition:
gateway $A$ can only succeed if $N_A$ is not $Q_A$-smooth.
This is equivalent to asking that $N_A$ has at least one prime factor
$q$ with $\leg{q}{A} = -1$.
\end{remark}

% ============================================================
\section{Forced Small Factors Are Inert}\label{sec:forced}
% ============================================================

For Case~B QR7 primes ($p \equiv 1 \pmod{24}$),
$N_A = (p+A)/4$ has forced small prime factors determined by $A \bmod 24$.
Table~\ref{tab:forced} records these and their Legendre symbols.

\begin{table}[ht]
\centering
\caption{Forced small prime factors of $N_A$ for $p \equiv 1 \pmod{24}$,
and whether they are QRs mod $A$.
``Yes'' means the forced factor lies in $Q_A$ and does not help reach $t_A$.}
\label{tab:forced}
\begin{tabular}{@{}ccccc@{}}
\toprule
$A$ & $N_A \bmod 6$ & Forced factor & $({\cdot}/A)$ & In $Q_A$? \\
\midrule
$7$  & $2$ & $2$ & $+1$ & \textbf{Yes} \\
$11$ & $3$ & $3$ & $+1$ & \textbf{Yes} \\
$19$ & $5$ & none & --- & --- \\
$23$ & $0$ & $2,3$ & $+1,+1$ & \textbf{Yes} \\
$31$ & $2$ & $2$ & $+1$ & \textbf{Yes} \\
$43$ & $5$ & none & --- & --- \\
$47$ & $0$ & $2,3$ & $+1,+1$ & \textbf{Yes} \\
$59$ & $3$ & $3$ & $+1$ & \textbf{Yes} \\
$71$ & $0$ & $2,3$ & $+1,+1$ & \textbf{Yes} \\
\bottomrule
\end{tabular}
\end{table}

\begin{lemma}[Forced factors are QR]\label{lem:forced-qr}
For all primes $A \equiv 3 \pmod{4}$ in the candidate set,
the forced small prime factors of $N_A$ are quadratic residues mod $A$.
\end{lemma}

\begin{proof}
For the factor $2$: $(2/A) = (-1)^{(A^2-1)/8}$.
For $A \equiv 7 \pmod{8}$: $(A^2-1)/8 = (A-1)(A+1)/8$;
since $8 \mid (A+1)$, this is even, giving $(2/A) = +1$.

For the factor $3$: $(3/A) = +1$ for $A \equiv 11 \pmod{12}$
(by quadratic reciprocity and direct check for each $A$ in the table).
\end{proof}

\begin{corollary}\label{cor:forced-inert}
The forced small prime factors of $N_A$ contribute only QR elements
to $S_A$ and cannot produce the NQR target $t_A$.
The gateway can succeed only via \emph{unforced} prime factors of $N_A$
that are NQR mod $A$.
\end{corollary}

This is structurally telling: the ``free gifts'' built into the framework
(the guaranteed divisibilities) are precisely the primes that are inert
in the relevant quadratic extension $\mathbb{Q}(\sqrt{A^*})$.

% ============================================================
\section{Character Sum Criterion}\label{sec:charsum}
% ============================================================

The count of divisors of $N_A^2$ in the target class is:
\begin{equation}\label{eq:count}
  \mathcal{N}_A(p) \;=\; \frac{1}{A-1}\sum_{\chi \bmod A}\chi(t_A)^{-1}
  \sum_{d \mid N_A^2}\chi(d).
\end{equation}

\begin{lemma}[Character sum factorisation]\label{lem:char-factor}
For $\chi \bmod A$ and $N$ with $\gcd(N,A)=1$:
\[
  \sum_{d \mid N^2}\chi(d)
  = \prod_{q^a \| N}\sum_{k=0}^{2a}\chi(q)^k.
\]
For the quadratic character $\chi_2$ (Legendre symbol):
\[
  \sum_{d \mid N^2}\chi_2(d) \;=\; \tau(M^2),
\]
where $M$ is the $Q_A$-smooth part of $N$ (product of QR prime factors with multiplicity).
\end{lemma}

\begin{proof}
The first identity is multiplicativity.
For $\chi_2$ and $(q/A) = -1$: $\sum_{k=0}^{2a}(-1)^k = 1$ (since $2a$ is even).
For $(q/A) = +1$: $\sum_{k=0}^{2a} 1 = 2a+1 = \tau(q^{2a})$.
The product over QR factors gives $\tau(M^2)$; NQR factors contribute $1$ each.
\end{proof}

From~\eqref{eq:count} and Lemma~\ref{lem:char-factor}:
\begin{equation}\label{eq:main-error}
  \mathcal{N}_A(p) \;=\; \frac{\tau(N_A^2) - \tau(M_A^2)}{A-1} + E_A(p),
\end{equation}
where $E_A(p)$ involves only characters of order $> 2$.
The main term $\tau(N_A^2) - \tau(M_A^2)$ is positive iff $N_A$ is not $Q_A$-smooth.

\textbf{Key consequence:}
A sufficient condition for $\mathcal{N}_A(p) \ge 1$ is:
\[
  \tau(N_A^2) - \tau(M_A^2) \;\ge\; (A-1)(1 + |E_A(p)|).
\]
When $N_A$ has many NQR prime factors (so $\tau(N_A^2) \gg \tau(M_A^2)$)
and $A$ is small, this is easily satisfied.
The difficult case is when $N_A$ has exactly one NQR prime factor of low
valuation---then the main term is small and the error may dominate.

% ============================================================
\section{The Multi-$A$ Strategy}\label{sec:multiA}
% ============================================================

Let $\mathcal{A} = \{A_1, \ldots, A_K\}$ be the set of $K$ candidate
gateway parameters.  Gateway $\mathcal{A}$ \emph{completely fails} for $p$
if $N_{A_i}$ is $Q_{A_i}$-smooth for all $i$.

\begin{theorem}[Multi-$A$ failure count]\label{thm:multiA}
Treating the $Q_{A_i}$-smoothness conditions as asymptotically independent
(justified below):
\[
  |\{p \le X : \text{all gateways in } \mathcal{A} \text{ fail}\}|
  \;\lesssim\; C \cdot \frac{X}{(\log X)^{1 + K/2}},
\]
where the exponent $K/2$ reflects the $(\log X)^{-1/2}$ density
of $Q_A$-smooth integers (Lemma~\ref{lem:qa-density}).
For $K \ge 3$, this count is $o(X/(\log X))$---the failures
have density zero among all primes.
\end{theorem}

\begin{lemma}[Density of $Q_A$-smooth integers]\label{lem:qa-density}
$|\{n \le X : n \text{ is } Q_A\text{-smooth}\}| \sim C_A X / (\log X)^{1/2}$
for a positive constant $C_A$.
\end{lemma}

\begin{proof}
Standard Selberg sieve argument: primes split equally into QR and NQR classes
mod $A$ by Dirichlet's theorem; the generating Dirichlet series has a factor
$\prod_p (1 - 1/p^2)^{-1/2}$ relative to the zeta function, giving the
stated $(\log X)^{-1/2}$ factor.
See \cite{Tenenbaum1984, Halberstam1974}.
\end{proof}

\textbf{Independence justification.}
The integers $N_{A_i} = (p + A_i)/4$ for distinct $i$ differ by $(A_i - A_j)/4$,
a fixed bounded quantity.  Their GCD is at most $(A_i - A_j)/4 = O(1)$,
so they share only $O(1)$ prime factors.
The conditions ``$N_{A_i}$ is $Q_{A_i}$-smooth'' involve different moduli $A_i$
applied to different (coprime) integers $N_{A_i}$;
standard correlation estimates (Elliott~\cite{Elliott1980})
show these events are asymptotically independent over primes $p$.

% ============================================================
\section{A Conditional Theorem}\label{sec:conditional}
% ============================================================

\begin{conjecture}[Joint $Q_A$-smoothness equidistribution]\label{conj:equidist}
Let $\mathcal{A} = \{A_1, \ldots, A_K\}$, each $A_i \equiv 3 \pmod{4}$ prime.
Then:
\[
  \#\bigl\{p \le X : N_{A_i} \text{ is } Q_{A_i}\text{-smooth},\; \forall i \le K\bigr\}
  = O\!\left(\frac{X}{(\log X)^{1 + \varepsilon}}\right)
\]
for some explicit $\varepsilon = \varepsilon(K, \mathcal{A}) > 0$.
\end{conjecture}

\begin{remark}
For $K = 1$: provable unconditionally with $\varepsilon = 1/2$ via the Selberg sieve.
For $K \ge 2$: implied by the Elliott--Halberstam conjecture
and consistent with the Generalised Riemann Hypothesis.
\end{remark}

\begin{theorem}[Conditional closure]\label{thm:conditional}
Assume Conjecture~\ref{conj:equidist} for the 24-element candidate set
$\mathcal{A}$ of~\cite{Gateway2026}.
Then the set of Case~B QR7 primes for which every gateway in $\mathcal{A}$ fails
has natural density zero, and for an explicitly computable $P_0 = P_0(\varepsilon)$,
the Erd\H{o}s--Straus conjecture holds for all $p > P_0$.
\end{theorem}

% ============================================================
\section{The Precise Analytic Gap}\label{sec:gap}
% ============================================================

The unconditional proof reduces to:

\begin{theorem}[Needed]\label{thm:needed}
For $K \ge 2$ fixed primes $A_1, \ldots, A_K \equiv 3 \pmod{4}$,
there exists explicit $\delta = \delta(K, \mathcal{A}) > 0$ such that
\[
  \#\bigl\{p \le X : N_{A_i} \text{ is } Q_{A_i}\text{-smooth},\;\forall i\bigr\}
  = O\!\left(\frac{X}{(\log X)^{1+\delta}}\right).
\]
\end{theorem}

This is a \emph{joint smoothness-in-subgroup} problem for prime shifts.
It sits at the intersection of:
\begin{enumerate}
  \item Hooley's $\Delta$-function and divisors in progressions~\cite{Hooley1979};
  \item Tenenbaum's anatomy-of-integers programme~\cite{Tenenbaum1984};
  \item The bilinear large sieve (Gallagher, Montgomery--Vaughan);
  \item Elliott's independence of multiplicative functions on shifts~\cite{Elliott1980}.
\end{enumerate}

For $K = 1$: Theorem~\ref{thm:needed} holds with $\delta = 1/2$ (Selberg sieve).
For $K \ge 2$: the joint bound requires controlling the correlation
$\text{Cov}(B_{A_1}, B_{A_2})$ where $B_{A_i}$ is the event
``$N_{A_i}$ is $Q_{A_i}$-smooth.''
The key difficulty: $N_{A_1}$ and $N_{A_2}$ are not independent
(both derived from the same prime $p$),
so a bilinear character sum estimate is needed.

\textbf{Most likely unconditional route:}
Apply the Bombieri--Vinogradov theorem for characters modulo $A_1 A_2 \cdots A_K$
to show that the Frobenius of $p$ in $\mathrm{Gal}(\mathbb{Q}(\zeta_{A_1 \cdots A_K})/\mathbb{Q})$
acts non-trivially on some $Q_{A_i}$-coset for most primes $p$.
This typically gives the ``almost all'' conclusion.
An ``all'' conclusion (required for a full proof) would need Maier-matrix or
exceptional-set methods (Pintz~\cite{Pintz1997}).

% ============================================================
\section{Nine Questions for Further Work}\label{sec:questions}
% ============================================================

\begin{question}[Galois group of simultaneous QR conditions]\label{q:galois}
Let $F$ be the compositum of quadratic fields
$\mathbb{Q}(\sqrt{A_i^*})$ for $i = 1, \ldots, K$
(where $A_i^* = (-1)^{(A_i-1)/2} A_i$).
The $Q_{A_i}$-smoothness conditions are: the Frobenius of each prime factor of
$N_{A_i}$ at $A_i$ is trivial in $\mathrm{Gal}(\mathbb{Q}(\sqrt{A_i^*})/\mathbb{Q})$.
Can the Chebotarev density theorem over $F$ (with explicit error)
quantify the joint failure probability?
Does $F$ have a natural interpretation in Egyptian fraction theory?
\end{question}

\begin{question}[Large sieve exponent for joint failure]\label{q:largesieve}
For $K = 2$, what is the sharpest achievable exponent $\delta$
in Theorem~\ref{thm:needed} via the large sieve inequality?
Is $\delta = 1$ achievable (i.e., can the count of doubly-failing primes
be bounded by $O(X / \log^2 X)$)?
\end{question}

\begin{question}[Anatomy of the hardest prime]\label{q:anatomy}
For $p = 32{,}349{,}601$ (requiring $A = 239$),
compute the factorizations of $(p + A_i)/4$ for all $A_i \in \{7, 11, 19, \ldots, 199\}$
and verify each is $Q_{A_i}$-smooth.
Is there an arithmetic pattern (e.g.\ unusual factorisation mod $\prod A_i$)
explaining this prime's resistance?
\end{question}

\begin{question}[$v$-adic depth and the failure threshold]\label{q:padic}
Define the \emph{NQR depth} of $(p, A)$ as
$\delta_A(p) = \min\{v_q(N_A) : q \mid N_A,\; (q/A) = -1\}$
($+\infty$ if $Q_A$-smooth).
When $\delta_A(p) \ge 2$, does the character sum
$\sum_{d \mid N_A^2} \chi_{\text{higher}}(d)$ become dominated by the main term,
unconditionally guaranteeing success?
Is $\delta_A(p) \ge 2$ alone sufficient for the gateway to succeed for $A \ge 11$?
\end{question}

\begin{question}[Class field interpretation]\label{q:classfield}
The Case~B condition (all prime factors of $(p+3)/4$ congruent to $1 \pmod 3$)
is equivalent to $(p+3)/4$ being a norm from $\mathbb{Z}[\omega]$
(the Eisenstein integers), up to 2-adic factors.
Can simultaneous gateway failure be expressed as a norm condition
in a ring of integers of a biquadratic field
$\mathbb{Q}(\sqrt{-3}, \sqrt{-7})$ or similar?
\end{question}

\begin{question}[Stabilisation of $\max A$]\label{q:maxA}
Let $A^*(X)$ be the maximum gateway parameter required for primes $\le X$.
Is $\lim_{X \to \infty} A^*(X)$ finite?
Can the Bombieri--Vinogradov theorem give $A^*(X) = O((\log X)^2)$,
analogously to the bound on the least quadratic non-residue?
\end{question}

\begin{question}[Elliott's programme and independence]\label{q:multiplicative}
Elliott's work on multiplicative functions on shifted primes asks when
$f_1(n + a)$ and $f_2(n + b)$ are asymptotically independent
for multiplicative $f_1, f_2$.
Does the recent Frantzikinakis--Host theorem
(or Klurman--Mangerel on multiplicative correlations)
apply to the pair $(r_{A_1}(N_{A_1}), r_{A_2}(N_{A_2}))$
where $r_{A_i}$ is the indicator of $Q_{A_i}$-smooth integers?
\end{question}

\begin{question}[The $k/p$ generalisation]\label{q:generalise}
For the generalised problem $k/p = 1/x + 1/y + 1/z$ (any $k \ge 4$),
the gateway target is $t_{A,k} = -k N_{A,k}^2 \bmod A$.
Lemma~\ref{lem:nqr-target} extends: $t_{A,k}$ is NQR iff $(-k/A) = -1$,
i.e., $k$ is a QR mod $A$.
For which $k$ does the bounded-$A$ phenomenon persist?
Is the multi-$A$ strategy equally effective for $k = 5, 6, 7$?
\end{question}

\begin{question}[Computation to $10^{10}$]\label{q:compute}
Does any prime $p \in (10^9, 10^{10})$ require $A > 239$?
If $A^*(10^{10}) = 239$ (no new maximum), this strongly supports
the bounded-$A$ conjecture and, combined with Theorem~\ref{thm:conditional},
makes an explicit threshold $P_0$ computationally accessible.
If $A^*(10^{10}) > 239$, the new maximum prime provides a case study
analogous to $p = 32{,}349{,}601$ for further analysis.
\end{question}

% ============================================================
\section{Conclusion}
% ============================================================

The central results are:
\begin{itemize}
\item \textbf{Lemma~\ref{lem:nqr-target}}: the gateway target is always
  a NQR mod $A$ (a clean one-line consequence of $(-1/A) = -1$).
\item \textbf{Theorem~\ref{thm:failure-char}}: gateway $A$ fails iff
  $N_A$ is $Q_A$-smooth.
  This replaces the vague ``equidistribution'' condition in the companion
  paper with a precise algebraic obstruction.
\item \textbf{Theorem~\ref{thm:conditional}}: the Erd\H{o}s--Straus conjecture
  follows from a standard analytic hypothesis
  (Conjecture~\ref{conj:equidist}) that is weaker than Elliott--Halberstam.
\end{itemize}

The remaining unconditional gap is the joint $Q_A$-smoothness bound
(Theorem~\ref{thm:needed}), which appears tractable via bilinear
character sums and the large sieve but requires new estimates for
correlated prime-shifted integers.

\paragraph{Code.}
\url{https://github.com/m4cd4r4/erdos-straus-gateway}

\begin{thebibliography}{99}

\bibitem{Gateway2026}
M.~\'O~Murch\'u,
\emph{A Finite Algebraic Covering System for the
Erd\H{o}s--Straus Conjecture to $10^9$}, preprint, 2026.

\bibitem{Elliott1980}
P.~D.~T.~A.~Elliott,
\emph{Probabilistic Number Theory, I}, Springer, 1979.

\bibitem{Halberstam1974}
H.~Halberstam and H.-E.~Richert,
\emph{Sieve Methods}, Academic Press, 1974.

\bibitem{Hooley1979}
C.~Hooley,
\emph{On a new technique and its applications to the theory of numbers},
Proc.\ London Math.\ Soc.\ (3) \textbf{38} (1979), 115--151.

\bibitem{LMO1979}
J.~C.~Lagarias, H.~L.~Montgomery, and A.~M.~Odlyzko,
\emph{A bound for the least prime ideal in the Chebotarev density theorem},
Invent.\ Math.\ \textbf{54} (1979), 271--296.

\bibitem{Pintz1997}
J.~Pintz,
\emph{Very large gaps between consecutive primes},
J.\ Number Theory \textbf{63} (1997), 286--301.

\bibitem{Tenenbaum1984}
G.~Tenenbaum,
\emph{Sur la probabilit\'e qu'un entier poss\`ede un diviseur dans un
intervalle donn\'e}, Compositio Math.\ \textbf{51} (1984), 243--263.

\end{thebibliography}

\end{document}
